\begin{center}\textit{By L. Cullen (triponacci)}\end{center}


%==========================================================
\section*{Coset Index and Gauge Confinement}
%==========================================================

The arithmetic of prime splitting extends directly to the physical gauge hierarchy via the Coset Index (CI), defined as the index of the cyclic subgroup generated by the golden ratio $\vp$:
\begin{equation}
\mathrm{CI} = \frac{p-1}{\mathrm{ord}(\vp, p)}
\end{equation}
The Standard Model gauge channels SU(3), U(1), and SU(2) correspond to the prime tower CI values $\{2, 3, 4\}$. Dynamically, the CI determines the error scaling parameter $\Upsilon = 1/\mathrm{CI}$ in the continuous error ODE. Only the SU(3) channel ($\mathrm{CI} = 2 \implies \Upsilon = 0.5$) is \textbf{supercritical} ($\Upsilon > 1/\vp^2 \approx 0.382$), yielding an error growth root $r_1 \approx 1.781 > \vp$. This strict supercriticality proves that strong-force confinement is an information-theoretic necessity; without confinement, the non-associative topological friction would shatter the container.

\begin{thebibliography}{9}

\bibitem{neukirch}
J.~Neukirch,
\emph{Algebraic Number Theory},
Grundlehren der mathematischen Wissenschaften, vol.~322, Springer, 1999.

\bibitem{milne_cft}
J.~S.~Milne,
``Class Field Theory,''
lecture notes, v4.03, 2020.

\bibitem{serre_local}
J.-P.~Serre,
\emph{Local Fields},
Graduate Texts in Mathematics, vol.~67, Springer, 1979.

\bibitem{wall1960}
D.~D.~Wall,
``Fibonacci Series Modulo $m$,''
\textit{American Mathematical Monthly}, vol.~67, no.~6, pp.~525--532, 1960.

\bibitem{cassels-frohlich}
J.~W.~S.~Cassels and A.~Fr\"ohlich (eds.),
\emph{Algebraic Number Theory},
Academic Press, 1967.

\bibitem{lidl-niederreiter}
R.~Lidl and H.~Niederreiter,
\emph{Finite Fields}, 2nd ed.,
Cambridge University Press, 1997.

\bibitem{scholze2012}
P.~Scholze,
``Perfectoid Spaces,''
\textit{Publ.\ Math.\ Inst.\ Hautes \'Etudes Sci.}, vol.~116, pp.~245--313, 2012.

\end{thebibliography}

